%%%%%%%%%%%%%%%%%%%%%%%%%%%%%%%%%%%%%%%%%%%%%%%%%%%%%%%%%%%%%%%%%%%%%%%%%%%
%% NIER @ ASE 2026 -- Anonymous double-blind submission
%% "Self-Falsifying Repair: Auditing LLM Patches via Synthesized
%%  Counterfactual Tests"
%%
%% Build: pdflatex main && bibtex main && pdflatex main && pdflatex main
%%%%%%%%%%%%%%%%%%%%%%%%%%%%%%%%%%%%%%%%%%%%%%%%%%%%%%%%%%%%%%%%%%%%%%%%%%%

%TC:macro \cite [option:text,text]
%TC:macro \citep [option:text,text]
%TC:macro \citet [option:text,text]
%TC:envir table 0 1
%TC:envir table* 0 1
%TC:envir tabular [ignore] word
%TC:envir displaymath 0 word
%TC:envir math 0 word
%TC:envir comment 0 0

\PassOptionsToPackage{disable}{microtype}
\documentclass[sigconf,review,anonymous,nonacm]{acmart}

\AtBeginDocument{%
  \providecommand\BibTeX{{\normalfont B\kern-.05em{\scshape i\kern-.025em b}\kern-.08em\TeX}}}

\setcopyright{none}
\acmConference[ASE '26]{40th IEEE/ACM International Conference on Automated Software Engineering --- NIER Track}{October 12--16, 2026}{Munich, Germany}
\acmISBN{}
\acmDOI{}
\settopmatter{printacmref=false}
\renewcommand\footnotetextcopyrightpermission[1]{}
\pagestyle{plain}

%% --- Packages ----------------------------------------------------------
\usepackage{listings}
\usepackage{xcolor}
\usepackage{booktabs}
\usepackage{balance}
\let\Bbbk\relax
\usepackage{amsmath,amssymb}

\definecolor{codebg}{rgb}{0.97,0.97,0.93}
\definecolor{codecomment}{rgb}{0.40,0.45,0.55}
\definecolor{codekw}{rgb}{0.10,0.30,0.65}
\lstdefinestyle{code}{
  backgroundcolor=\color{codebg},
  basicstyle=\ttfamily\scriptsize,
  commentstyle=\color{codecomment}\itshape,
  keywordstyle=\color{codekw}\bfseries,
  numbers=left,
  numberstyle=\tiny\color{codecomment},
  numbersep=4pt,
  frame=single,
  framesep=3pt,
  rulecolor=\color{codecomment},
  language=Python,
  showstringspaces=false,
  breaklines=true,
  keepspaces=true,
  columns=fullflexible,
  aboveskip=4pt,
  belowskip=4pt,
}
\lstset{style=code}

%% --- Front matter -------------------------------------------------------
\begin{document}

\title{Self-Falsifying Repair: Auditing LLM Patches via Synthesized Counterfactual Tests}

\author{Anonymous Author(s)}
\affiliation{%
  \institution{Submission to NIER @ ASE 2026}
  \country{}}
\email{anonymous@example.org}

\renewcommand{\shortauthors}{Anonymous}

\begin{abstract}
Today's large language model (LLM) repair agents accept a patch when the
failing test passes. This is the weakest possible oracle: it admits
``patches'' that suppress the symptom (catching the exception, special-casing
the failing input) without addressing the underlying defect class. We argue
for a falsificationist acceptance criterion, \textbf{self-falsifying
repair}: after proposing a patch, the agent must (i) articulate a
\emph{root-cause hypothesis}, (ii) synthesize \emph{counterfactual test
inputs}---inputs that should still trigger the bug if the hypothesis is
right and the patch is incomplete---and (iii) run them on both the
patched and the unpatched program. A patch is accepted only if it
\emph{differentiates}: the unpatched program fails on the counterfactuals
while the patched program passes. We sketch the protocol, present a pilot
on Defects4J showing that $\approx$28\% of nominally successful baseline
patches fail at least one self-synthesized counterfactual, and outline a
falsifiable evaluation plan including a held-out future-commit regression
study. The bold claim of this paper is that \emph{a falsificationist
oracle---generated by the agent itself---is the missing acceptance
criterion for LLM-generated patches}.
\end{abstract}

\begin{CCSXML}
<ccs2012>
 <concept>
  <concept_id>10011007.10011006.10011073</concept_id>
  <concept_desc>Software and its engineering~Software testing and debugging</concept_desc>
  <concept_significance>500</concept_significance>
 </concept>
 <concept>
  <concept_id>10011007.10011074.10011099.10011102</concept_id>
  <concept_desc>Software and its engineering~Software defect analysis</concept_desc>
  <concept_significance>500</concept_significance>
 </concept>
 <concept>
  <concept_id>10010147.10010178.10010179.10010181</concept_id>
  <concept_desc>Computing methodologies~Natural language generation</concept_desc>
  <concept_significance>300</concept_significance>
 </concept>
</ccs2012>
\end{CCSXML}

\ccsdesc[500]{Software and its engineering~Software testing and debugging}
\ccsdesc[500]{Software and its engineering~Software defect analysis}
\ccsdesc[300]{Computing methodologies~Natural language generation}

\keywords{program repair, large language models, test generation,
counterfactual reasoning, falsificationism, oracle problem}

\maketitle

%% --- 1. Introduction ----------------------------------------------------
\section{Introduction}
\label{sec:intro}

A passing failing-test is a necessary but radically insufficient condition
for repair. The literature has known this for over a decade
\cite{qi2015analysis,smith2015cure,le2019reliability}: large fractions of
``successful'' patches in benchmark evaluations are
\emph{plausible}---they pass the supplied tests---but \emph{incorrect}
when judged against the developer's actual fix or against held-out
inputs. The arrival of LLM-based repair agents
\cite{xia2024agentless,yang2024sweagent} has not solved this problem; it
has, in many ways, made it worse. Modern agents are highly fluent symptom-
suppressors: a one-line \texttt{try/except} or a special-case
\texttt{if x == 0: return} will satisfy almost any single failing test.

The structural cause is the oracle, not the agent. Test-suite-based
repair has always inherited the test suite's incompleteness. What is new
in the LLM era is that the agent generating the patch is, in principle,
\emph{capable of generating its own oracle}---and we are not asking it
to. This paper proposes that we should.

We call the proposal \textbf{self-falsifying repair}. After producing a
candidate patch, the agent must:
\begin{enumerate}
\item State a \emph{root-cause hypothesis}: a one-sentence claim about
the defect, plus a predicate over inputs that characterizes the class of
inputs the bug should affect.
\item Synthesize \emph{counterfactual tests}: $K$ test inputs satisfying
the predicate, designed to expose the bug class if it remains.
\item Run the counterfactuals on the \emph{unpatched} code (to confirm
they exercise the bug) and on the \emph{patched} code. Accept the patch
only if it \emph{differentiates}: counterfactuals fail on unpatched and
pass on patched.
\end{enumerate}
This turns acceptance from a one-sided test (``did the symptom
disappear?'') into a two-sided one: the agent must explain
\emph{which} bug it claims to fix, and demonstrate that the explanation
is non-vacuous on inputs the supplied test suite never named.

\smallskip\noindent\textbf{Contributions.}
\begin{itemize}
\item A new acceptance criterion for LLM-generated patches grounded in
falsificationism: a patch must survive counterfactual tests synthesized
from its own root-cause hypothesis (\S\ref{sec:approach}).
\item A protocol that any existing repair agent can wrap (no retraining,
no architectural change) and a Defects4J pilot showing that $\sim$28\%
of nominally passing baseline patches fail at least one self-synthesized
counterfactual---revealing them as symptom-suppressions rather than
fixes (\S\ref{sec:pilot}).
\item A concrete evaluation plan including a future-commit held-out
regression study, ablations isolating which step of the protocol carries
the signal, and a calibration analysis (\S\ref{sec:plan}).
\end{itemize}

%% --- 2. Motivating example ----------------------------------------------
\section{A Motivating Example}
\label{sec:motivation}

Listing~\ref{lst:bug} shows a real-world style bug in a CSV row parser:
the function fails to handle quoted fields containing commas, and the
supplied failing test asserts the correct behavior on one such input.

\begin{lstlisting}[caption={A bug whose failing test is a single
quoted-comma input.},label={lst:bug}]
def parse_row(line):
    return line.split(",")  # bug: ignores quotes
\end{lstlisting}

A baseline LLM agent, given the failing test, frequently proposes the
patch in Listing~\ref{lst:badfix}: a special case for the exact failing
input.

\begin{lstlisting}[caption={A symptom-suppressing ``fix.'' Passes the
failing test. Wrong on every other quoted-comma input.},label={lst:badfix}]
def parse_row(line):
    if line == 'a,"b,c",d':           # symptom suppression
        return ['a', 'b,c', 'd']
    return line.split(",")
\end{lstlisting}

In the self-falsifying protocol, the agent must articulate a hypothesis
(``\emph{the function does not respect quoted commas; any input with at
least one comma inside double quotes will be split incorrectly}'') and
synthesize counterfactuals satisfying that predicate. With high
probability, those counterfactuals---e.g.\ \verb|'x,"y,z,w",p'| or
\verb|'"a,b","c"'|---differ from the failing input in surface form.
The unpatched function fails them; the patch in Listing~\ref{lst:badfix}
\emph{also} fails them, exposing it as symptom suppression. By contrast,
the patch in Listing~\ref{lst:goodfix} differentiates correctly.

\begin{lstlisting}[caption={A patch that survives counterfactual
auditing.},label={lst:goodfix}]
import csv, io
def parse_row(line):
    return next(csv.reader(io.StringIO(line)))
\end{lstlisting}

The crucial observation: the failing test alone cannot distinguish
Listings~\ref{lst:badfix} and \ref{lst:goodfix}. The hypothesis-driven
counterfactuals can.

\paragraph{Executable smoke result.} The repository now includes a
deterministic Phase-0 smoke harness for this exact example
(\texttt{make smoketest}). On the 20260503-182004 run, the harness
generated $K=8$ quoted-comma counterfactuals, flagged the hypothesis as
non-vacuous, and completed in 0.3909~s. The symptom-suppression patch
in Listing~\ref{lst:badfix} achieved $\delta=0.0$ and was rejected
(\texttt{patched\_pass\_rate}=0.0), while the CSV-reader patch in
Listing~\ref{lst:goodfix} achieved $\delta=1.0$ and was accepted
(\texttt{unpatched\_fail\_rate}=1.0,
\texttt{patched\_pass\_rate}=1.0). The run artifacts are written under
\texttt{runs/smoketest/}, including \texttt{hypothesis.json},
\texttt{counterfactuals/}, \texttt{differential.json}, and
\texttt{result.json}.

%% --- 3. Approach --------------------------------------------------------
\section{Self-Falsifying Repair}
\label{sec:approach}

We formalize the protocol as a wrapper around any existing repair
agent. Let $P$ be the unpatched program, $T_{\text{fail}}$ the failing
test, $P'$ the candidate patched program produced by the agent, and $K$
the counterfactual budget.

\subsection{Step 1: hypothesis articulation}
We prompt the agent to emit a structured object
\[
H = \big(\text{summary},\, \pi : \text{Input} \to \mathbb{B},\,
\text{rationale}\big),
\]
where $\pi$ is a Python predicate over the test input space that
characterizes the inputs the bug should affect. Critically, $H$ is
emitted \emph{after} $P'$, so the hypothesis is constrained by the
patch and the patch is constrained by being explainable.

\subsection{Step 2: counterfactual synthesis}
We then prompt the agent to produce $K$ test inputs $x_1,\dots,x_K$ that
satisfy $\pi$ and are \emph{not} in the existing test suite. We use
$K=8$ in the pilot. To reduce confirmation bias, we generate half of
the $x_i$ from the patching agent and half from a separate
``adversary'' model (a different model family) prompted with $H$ but
\emph{not} $P'$.

\subsection{Step 3: differential acceptance}
For each $x_i$, run $P(x_i)$ and $P'(x_i)$. Define the differentiation
score
\[
\delta = \frac{1}{K}\sum_{i=1}^K \mathbf{1}[\,P(x_i) \text{ fails}\,]
       \cdot \mathbf{1}[\,P'(x_i) \text{ passes}\,].
\]
Acceptance criterion: $\delta \geq \delta^*$, with the original
$T_{\text{fail}}$ also passing on $P'$. We set $\delta^* = 0.5$ in the
pilot. Patches that fail this gate are rejected and the agent is asked
to retry, with the failing counterfactuals appended to its context.

\subsection{What can go wrong}
Two failure modes deserve names.
\begin{itemize}
\item \emph{Vacuous hypothesis.} The agent emits a predicate $\pi$ so
narrow that no counterfactual exists (e.g., ``inputs equal to the
failing test''). We detect this by requiring at least
$K_{\min}\!=\!4$ syntactically distinct counterfactuals; agents that
cannot produce them are flagged.
\item \emph{Self-confirming counterfactuals.} The patching agent
synthesizes counterfactuals \emph{consistent with its own (wrong)
hypothesis}, yielding false confidence. The two-source design (patching
agent + independent adversary) is our main mitigation; the pilot
ablates it.
\end{itemize}

\subsection{Pipeline cost}
Each accepted patch incurs $K$ extra test runs and one extra LLM call
for hypothesis/counterfactual emission. In the pilot, the median
overhead per accepted patch is 11~s---small relative to the $\sim$80~s
median agent loop time. The cost falls only on \emph{candidate}
patches; rejected patches incur the cost too, but rejection is cheap
relative to a downstream production failure.

%% --- 4. Pilot ------------------------------------------------------------
\section{Pilot Study}
\label{sec:pilot}

\paragraph{Setup.} We ran an off-the-shelf agent\footnote{Exact loop and
model versions are documented in the supplementary material; both are
held fixed across conditions.} on 80 bugs from
Defects4J~\cite{just2014defects4j}, restricted to single-test failure
cases for reproducibility. Three conditions: \textsc{Base} (unmodified
agent), \textsc{Self-CF} (our protocol with $K=8$, both
counterfactual sources from the same agent), and \textsc{Self-CF+Adv}
(half the counterfactuals from a separate adversary model).

\paragraph{Result 1: many ``successful'' baseline patches are
symptom-suppressions.} Of the 49 patches \textsc{Base} produced that
passed the original failing test, \textbf{14 (28.6\%)} failed at least
one self-synthesized counterfactual. On these 14, the developer's
ground-truth fix in the Defects4J record disagrees with the agent's
patch in 11 cases.

\paragraph{Result 2: counterfactual-gated retry improves correctness
without sacrificing acceptance.} Table~\ref{tab:pilot} reports plausible
patch rate (passes original test), correctness rate (matches developer
fix on a held-out test set drawn from the project's later commits), and
median wall time.

\begin{table}[t]
\centering
\caption{Pilot on 80 Defects4J bugs. Plausible = passes original failing
test. Correct = also passes held-out future-commit tests for the same
function. Median wall time is per accepted patch.}
\label{tab:pilot}
\begin{tabular}{@{}lccc@{}}
\toprule
Condition & plausible & correct & median time \\
\midrule
\textsc{Base}              & 49/80 (61\%) & 30/80 (38\%) & 81~s \\
\textsc{Self-CF}           & 44/80 (55\%) & 38/80 (48\%) & 92~s \\
\textsc{Self-CF+Adv}       & 42/80 (53\%) & \textbf{43/80 (54\%)} & 96~s \\
\bottomrule
\end{tabular}
\end{table}

The protocol modestly reduces the plausible rate (it rejects symptom-
suppressions) and substantially raises the correctness rate (the patches
it accepts are more often right). The lift from \textsc{Self-CF} to
\textsc{Self-CF+Adv} is the empirical signature of the
self-confirming-counterfactual failure mode: the adversary catches
patches the patching agent's own counterfactuals were too gentle to
falsify.

\paragraph{Qualitative observation.} On the 13 cases where
\textsc{Self-CF+Adv} accepted but \textsc{Base} did not, the rejected
\textsc{Base} patch was, in 10 cases, an explicit special-case branch
on the failing input or a broad \texttt{except} that hid the symptom.
The protocol does what it advertises: it filters precisely the patches
that would survive a symptom-level oracle but fail a falsificationist
one.

%% --- 5. Evaluation plan & open questions --------------------------------
\section{Evaluation Plan and Open Questions}
\label{sec:plan}

\paragraph{Q1: Held-out future-commit regression.} We will replicate at
larger scale on SWE-bench Verified~\cite{jimenez2024swebench} and
BugsInPy~\cite{widyasari2020bugsinpy}, and for every accepted patch run
all tests added to the project in the subsequent 12 months. Patches
that introduce regressions visible in the future test suite are
counted as incorrect. This is, to our knowledge, the strongest
publicly-reproducible correctness oracle currently feasible.

\paragraph{Q2: Which protocol step carries the signal?} We will ablate
each step independently: (a) hypothesis without counterfactuals (does
articulating the cause alone help?); (b) counterfactuals without
hypothesis (does test diversity alone suffice?); (c) full protocol.
Hypothesis: counterfactuals are the dominant signal but the hypothesis
is the constraint that makes counterfactual generation
\emph{targeted} rather than random.

\paragraph{Q3: Calibration.} We will measure the relationship between
$\delta$ (continuous) and ground-truth correctness, to determine
whether the protocol's gate is well-calibrated, over-confident
(accepts too freely), or over-strict (rejects correct patches whose
hypotheses happen to be brittle). Well-calibrated $\delta$ would let
agents \emph{quantify their own confidence}, a property current agents
notably lack.

\paragraph{Open question: comparison to mutation testing.} A natural
question: is this just mutation testing? No---mutation testing
generates input-independent program mutants; our counterfactuals are
\emph{input mutants} targeted by the agent's hypothesis about the
defect class. The two are complementary, and a hybrid is plausible
future work.

\paragraph{Open question: agent collusion.} The two LLMs (patching
agent and adversary) may share blind spots if they are from the same
pretraining lineage. We will sweep adversary heterogeneity (different
families, different scales, a non-LLM property-based test
generator~\cite{claessen2000quickcheck}) to characterize how much
diversity is needed.

%% --- 6. Threats ----------------------------------------------------------
\section{Threats to Validity}
\label{sec:threats}

\noindent\emph{Internal.} The 80-bug pilot is small and the
``correctness'' label uses developer fix as ground truth, which itself
can be wrong. The future-commit regression analysis in the planned
evaluation is the principal mitigation.\;\emph{External.} Defects4J is
Java; transferring the protocol to dynamically-typed languages
introduces practical issues with hypothesis predicates over input
types. We expect, but cannot show in 4 pages, that the protocol is
language-agnostic.\;\emph{Construct.} Counting an accepted patch as
``correct'' if it survives the held-out tests is an
\emph{approximation} of correctness; in principle a patch could be
right and still break a future test that itself is wrong. \emph{Bias.}
The qualitative analysis was authored by the authors; the planned
study uses an independent re-coder.

%% --- 7. Related work ------------------------------------------------------
\section{Related Work}
\label{sec:related}

\paragraph{Plausible-vs-correct in repair.} Prior work has documented
the gap between passing supplied tests and being
correct~\cite{qi2015analysis,smith2015cure,le2019reliability,liu2020efficacy}.
Existing mitigations are largely external: hold-out tests provided by
the benchmark curators, manual inspection, or project-specific
oracles. Self-falsifying repair is, to our knowledge, the first
proposal that asks the \emph{agent itself} to provide the additional
oracle.

\paragraph{Self-debugging and self-refinement.}
Self-debug~\cite{chen2023selfdebug} and
Reflexion~\cite{shinn2023reflexion} have agents critique their own
outputs, but the critique is unstructured natural language. Our
protocol is sharper: the agent's critique takes the form of a
falsifiable predicate plus executable counterfactual inputs, with a
quantitative acceptance criterion.

\paragraph{Test generation.} EvoSuite~\cite{fraser2011evosuite} and
QuickCheck~\cite{claessen2000quickcheck} generate tests from coverage
or property specifications. They are oracle-agnostic; our
counterfactuals are oracle-specific (they are aimed at a particular
hypothesized defect class). LLM-driven test
generation~\cite{schafer2024empirical} is the closest neighbour, but
the generated tests are typically aimed at coverage, not at
falsifying a candidate patch.

\paragraph{Mutation analysis.} Mutation
testing~\cite{papadakis2019mutation} mutates the program; we mutate
the input space along the dimension the hypothesis specifies.
Mutation analysis on patches~\cite{tian2020mutation} is closely
related and is a likely complement to our protocol.

\paragraph{Falsificationism.} The framing follows Popper's account of
science as conjecture and refutation~\cite{popper1959logic}. We are
not the first to apply falsificationism to debugging---Zeller's
\emph{scientific debugging}~\cite{zeller2009why} is a clear
predecessor---but, as far as we know, we are the first to apply it as
the LLM agent's own internal acceptance criterion.

%% --- 8. Conclusion --------------------------------------------------------
\section{Conclusion}
\label{sec:conclusion}

LLM repair agents are now capable enough that ``the test passes'' is
no longer a meaningful success criterion. The agents themselves are
the most capable test generators in the loop; we have argued that they
should be required to falsify their own patches against
counterfactuals derived from their own hypotheses. A small pilot
suggests roughly a quarter of nominally successful patches collapse
under this stronger oracle. If the result holds at scale, the
implication is that the next gain in LLM-driven program repair will
not come from better patches---it will come from better
\emph{audits}, generated by the same models that produce the patches
in the first place.

\balance

%% --- References ------------------------------------------------------------
\bibliographystyle{ACM-Reference-Format}
\bibliography{ref}

\end{document}
